\documentclass[../readings.tex]{subfiles}

\begin{document}

\subsection{Mathematical Notation and Proof Techniques}
\label{sub:math-notation}

Standard vocabulary for proofs and algorithms.

\subsubsection{Logic and Set Notation}

\addDef{Logical Quantifiers}{%
\begin{itemize}
  \item \textbf{Universal ($\forall$):}\enspace ``For all.''
  \item \textbf{Existential ($\exists$):}\enspace ``There exists.''
  \item \textbf{Example:}\enspace ``$\forall \varepsilon > 0,\;\exists \delta > 0$'' reads as ``for every positive $\varepsilon$ there is a positive $\delta$.''
\end{itemize}
}
\label{def:quantifiers}

\addDef{Implication and Equivalence}{%
\begin{itemize}
  \item \textbf{Implication ($P \Rightarrow Q$):}\enspace $P$ implies $Q$.
  \item \textbf{Equivalence ($P \Leftrightarrow Q$):}\enspace $P$ if and only if $Q$ (iff).
  \item \textbf{Contrapositive:}\enspace $\neg Q \Rightarrow \neg P$. Logically equivalent to $P \Rightarrow Q$.
\end{itemize}
}
\label{def:implication}

\addDef{Set Notation}{%
\begin{itemize}
  \item $x \in S$: $x$ is an element of $S$.
  \item $S \subseteq T$: $S$ is a subset of $T$.
  \item $S \cup T$: Union; $S \cap T$: Intersection; $S \setminus T$: Set difference.
  \item $\{x \in S : P(x)\}$: Set builder notation.
  \item \textbf{Number Sets:}\enspace $\fZ$ (integers), $\fR$ (reals), $\fC$ (complex), $\fR^n$ (vectors), $\fR^{m \times n}$ (matrices).
\end{itemize}
}
\label{def:set-notation}

\addExcr{%
\begin{enumerate}
  \item Write ``every invertible matrix has a nonzero determinant'' using quantifiers and implication.
  \item State the contrapositive of: ``if $\bA$ is invertible, then $\bA\bx = \bzero \Rightarrow \bx = \bzero$.''
  \item Let $S = \{1, 2, 3, 4\}$ and $T = \{3, 4, 5, 6\}$. Compute $S \cup T, S \cap T, S \setminus T$.
\end{enumerate}
}

\subsubsection{Proof Techniques}

\addDef{Direct Proof}{%
Assume $P$ is true and derive $Q$ through a chain of valid logical steps.
}
\label{def:direct-proof}

\addDef{Proof by Contradiction}{%
Assume $\neg P$ and derive a logical contradiction ($R \wedge \neg R$). Proves $P$ must hold.
}
\label{def:proof-contradiction}

\addDef{Proof by Induction}{%
To prove $P(n)$ holds for all integers $n \geq n_0$:
\begin{enumerate}
  \item \textbf{Base case:}\enspace Verify $P(n_0)$ directly.
  \item \textbf{Inductive step:}\enspace Assume $P(k)$ and prove $P(k+1)$.
\end{enumerate}
}
\label{def:proof-induction}

\addNote{%
\textbf{Numerical Proofs:}\enspace Often validate that an approximation is **stable** or that an error term **converges**. Frequently involves showing $|f(x) - \hat{f}(x)| \leq \varepsilon$ for small $\varepsilon$.
}

\addExcr{%
\begin{enumerate}
  \item Prove: if $\bA$ and $\bB$ are symmetric, then $\bA + \bB$ is symmetric.
  \item Prove by contradiction: if $\bA\bx = \bzero \Rightarrow \bx = \bzero$, then $\bA$ has no zero eigenvalue.
  \item By induction, prove $\sum_{k=1}^{n} k = \frac{n(n+1)}{2}$.
\end{enumerate}
}

\subsubsection{Notation Conventions}

\addNote{%
\textbf{Scalars, Vectors, and Matrices.}\enspace
\begin{itemize}
  \item \textbf{Scalars:}\enspace Lowercase italic ($a, b, x$).
  \item \textbf{Vectors:}\enspace Lowercase bold ($\bx, \by, \ba$). Default is **column vector**.
  \item \textbf{Matrices:}\enspace Uppercase bold ($\bA, \bB, \bQ$).
  \item \textbf{Indexing:}\enspace 1-based (unlike Python). $x_i$ is $i$-th entry; $a_{ij}$ is $(i,j)$ entry.
\end{itemize}
}

\addNote{%
\textbf{Transpose and Inverse.}\enspace
\begin{itemize}
  \item $\bA^T$: Transpose.
  \item $\bA^{-1}$: Matrix inverse (square, nonsingular matrices).
  \item $\bA^{-T}$: Transpose of the inverse $(\bA^{-1})^T$.
  \item $\bA^k$: $k$-th power. $\bA^0 = \bI$.
\end{itemize}
}

\addNote{%
\textbf{Sums and Products.}\enspace
\[
  \sum_{i=1}^{n} a_i = a_1 + \dots + a_n, \qquad
  \prod_{i=1}^{n} a_i = a_1 \dots a_n.
\]
}

\addNote{%
\textbf{Functions.}\enspace
\begin{itemize}
  \item $f : \fR^n \to \fR^m$: Maps $n$-vectors to $m$-vectors.
  \item $f \circ g$: Function composition $(f \circ g)(\bx) = f(g(\bx))$.
\end{itemize}
}

\addNote{%
\textbf{Asymptotic Notation.}\enspace
Scaling of algorithms is described using **Big-O notation** ($O(n^2), O(n^3)$). Details in the Computational Complexity section.
}

\end{document}
